(1) 取场的时间因子\(e^{-i\omega t}\)，定义\(\dot a=(-i\omega_0-\gamma_i-\gamma_e)a+\sqrt{2\gamma_e}s_{\rm in}\)与\(s_{\rm out}=s_{\rm in}-\sqrt{2\gamma_e}a\)，使\(|a|^2\)为腔内储能。自由无驱动时\(d|a|^2/dt=-2(\gamma_i+\gamma_e)|a|^2\)，所以这里\(\gamma\)表示振幅衰减率，而非能量衰减率。
(2) 令\(\Delta=\omega-\omega_0\)、\(a(t)=a_\omega e^{-i\omega t}\)，则\(a_\omega=\sqrt{2\gamma_e}s_{\rm in}/(\gamma_i+\gamma_e-i\Delta)\)。由输入—输出式有\(t(\Delta)=s_{\rm out}/s_{\rm in}=[\gamma_i-\gamma_e-i\Delta]/[\gamma_i+\gamma_e-i\Delta]\)，故\(T(\Delta)=|t|^2=[(\gamma_i-\gamma_e)^2+\Delta^2]/[(\gamma_i+\gamma_e)^2+\Delta^2]\)。
(3) 共振时若\(\gamma_i=\gamma_e\)，则\(t(0)=0\)，总线直通端完全熄灭，这就是临界耦合。由\(1-T=4\gamma_i\gamma_e/[(\gamma_i+\gamma_e)^2+\Delta^2]\)可得共振被吸收的功率分数，临界耦合时恰为1；若\(\gamma_i=0\)则不可能有真正能量吸收，即使共振会产生强相位变化。
(4) 透射谱的缺口\(1-T\)是宽度参数\(\Gamma=\gamma_i+\gamma_e\)的Lorentz函数，在\(|\Delta|=\Gamma\)时降为峰值的一半，所以全宽\(\Delta\omega_{\rm FWHM}=2(\gamma_i+\gamma_e)\)。加载品质因数\(Q_L=\omega_0/[2(\gamma_i+\gamma_e)]\)；无内损\(\gamma_i\to0\)时功率透射\(T\to1\)对所有频率成立，但相位仍在共振附近急变，不能把全通率误判为没有腔内储能。
